\section{An elementary transfer from prime counts to signed cancellation}
\label{sec:bootstrap}
There is a second route to the boundary estimates. It uses only finite
summation identities after importing a quantitative prime number theorem.
Write $L(t)=1+\lfloor\log_2 t\rfloor$ for the binary length of an integer
$t\ge1$, and let $\Lambda(p^j)=\log p$ for a prime $p$ and $j\ge1$,
with value zero elsewhere. Fix any integer $A\ge1$. The required input is
\begin{equation}\label{eq:bootstrap-input}
    |E(t)|\le Kt/L(t)^A,\qquad E(t)=\sum_{m\le t}(\Lambda(m)-1).
\end{equation}
Logs here denote rational enclosure names; the multiplier $K$ and every
binary length are integers.

\paragraph{Obtaining the input from ProveIt.}
The theorem \texttt{thm:theta-modulus-established} in the repository's
\texttt{pnt\_li\_ipra\_2026-11-15/tex/theta\_bridge.tex} gives, with
$C=2^{240}$ and $\vartheta(t)=\sum_{p\le t}\log p$,
\[
  t\ge2^{C(r+1)^{10}}\quad\Longrightarrow\quad
  |\vartheta(t)-t|\le2^{-r}t.
\]
It is the upstream quantitative theorem used in this section.
For a fully specified instance of \eqref{eq:bootstrap-input} take
\[
 D_A=2(4A+4)^{A+1},\qquad
 K=(C+2)^{A+1}+4C^A(20A)^{10A}+D_A.
\]
Indeed set $L=L(t)$ and choose the largest $k$ with $Ck^{10}\le L-1$
by bounded search. For $k=0$ use
$|\vartheta(t)-t|/t\le1+\log t\le1+L$ and $L\le C$.
For $k\ge1$ the theorem bounds this ratio by $2^{1-k}$, while
$L\le C(k+1)^{10}$. The power-series inequality
$e^u\ge u^d/d!$ and $\log2\ge1/2$ give
$(k+1)^d 2^{-k}\le 2(2d)^d$ for $d\ge1$.
This bounds the $\vartheta$ contribution by the first two terms of $K$.
For each prime $p\le\sqrt t$, its powers of exponent at least two contribute
in total at most $\log t$. Thus
$0\le \sum_{m\le t}\Lambda(m)-\vartheta(t)\le\sqrt t\log t$.
Since $t\ge2^{L-1}$, the same exponential inequality gives
$L^{A+1}/\sqrt t\le D_A$, as required.

\paragraph{The finite identity.}
Let $e(m)=\Lambda(m)-1$. The operator $Df(n)=f(n)\log n$ is a derivation
for Dirichlet convolution, by $\log(ab)=\log a+\log b$. Write $\mathbf1$
for the constant-one function. Factorization gives
$D\mathbf1=\mathbf1*\Lambda$. Differentiate $\mu*\mathbf1=[n=1]$
and convolve on the right with $\mu$ to obtain $D\mu=-\mu*\Lambda$.
Summing this equality for $n\le x$ and using
$\sum_{dm\le x}\mu(d)=1$ gives for each positive integer $x$
\begin{equation}\label{eq:bootstrap-identity}
 M(x)\log x
 =\sum_{d\le x}\mu(d)\log(x/d)-1
       -\sum_{dm\le x}\mu(d)e(m).
\end{equation}
All sums on the right are finite.

\paragraph{Gain of a half-power in one step.}
Set $B_0=1$ and recursively take
\[
 B_{j+1}=16^{A+1}+64\cdot4^A(K+2)(B_j+D_A+1).
\]
We prove $|M(x)|\le B_j x/L(x)^{j/2}$ for $0\le j\le A$. The case
$j=0$ uses $|\mu|\le1$. Put $s=j/2$, $L=L(x)$, and suppose $j<A$.
For $L<16$ use the trivial bound and $B_{j+1}\ge16^{A+1}$. Otherwise
choose $U=2^{\lfloor\sqrt L\rfloor}$. In the double sum of
\eqref{eq:bootstrap-identity}, the part $m\le U$ is
$\sum_{m\le U}e(m)M(\lfloor x/m\rfloor)$. The binary length of
$\lfloor x/m\rfloor$ is at least $L-\lfloor\sqrt L\rfloor\ge L/2$.
A dyadic block $2^h\le m<2^{h+1}$ contributes at most $2(K+2)$
to $\sum |e(m)|/m$, since \eqref{eq:bootstrap-input} gives
$\sum_{m\le y}|e(m)|\le(K+2)y$. Thus the small-$m$ part has magnitude
at most
\[
       4(K+2) 2^s B_j x\, L^{1/2-s}.
\]
The remaining part is exactly
$\sum_{d\le x/U}\mu(d)(E(\lfloor x/d\rfloor)-E(U))$.
Each argument of $E$ is at least $U$, with binary length at least
$\sqrt L$. Since $\sum_{d\le x}1/d\le 1+\log x\le1+L$, its magnitude
is at most $3Kx L^{1-A/2}$.

Finally summation by parts rewrites the first sum on the right of
\eqref{eq:bootstrap-identity} as
$\sum_{d<x} M(d)\log(1+1/d)$. For $d<\sqrt x$ bound
$|M(d)|\le d$; for $d\ge\sqrt x$ use the inductive bound and
$L(d)\ge L/4$. The result is bounded by
$\sqrt x+4^s B_j x/L^s$.
Use $L^s/\sqrt x\le D_A$ and divide \eqref{eq:bootstrap-identity} by
$\log x\ge L/4$. As $s+1/2\le A/2$, the three estimates together
are at most $B_{j+1}x/L^{s+1/2}$, proving the step. All block endpoints
come from integer roots or bit lengths. The fractional powers in the
bounds for $M$ disappear on squaring nonnegative quantities. For
quantifier-free induction one can fix an arbitrary upper bound $X$ and
carry the Boolean conjunction of the integer estimates for $1\le x\le X$
at stage $j$; every use of the inductive hypothesis is within this list.

Taking $A=2b$ proves an effective $|M(x)|\le B_{2b}x/L(x)^b$ for
each positive integer $b$. Thus the repository's explicit prime estimate
supplies whatever fixed logarithmic saving is needed in ordinary
logarithmically weighted sums of $\mu(n)/n$.
